\section{From labels to an exact decision}
\label{sec:application}

We now apply Theorem~\ref{thm:labeling} to a scaled discounted map.
The labels need not locate its fixed point directly. Instead, we construct
two integer comparison vectors on opposite sides of that point. Their
coordinatewise distance is small, so a label excluding one vector from a
boundary strip also allows that boundary to move past a smaller strip
without excluding either vector. Repeating this operation gives the
accuracy required by Theorem~\ref{thm:discount}.

\subsection{Two nearby integer comparison vectors}

Use the integers $W,Q,A_0,H$ of Section~\ref{sec:discount}, and set
\begin{equation}
\begin{gathered}
 N=32H^2,\qquad \lambda=1/N=\lambda_0,\qquad \gamma=1-1/N,\\
 K=256(N+1),\qquad S=16HK,\qquad J=SW+N+1.
\end{gathered}
\label{eq:application-parameters}
\end{equation}
Let
\[
 F(t)=S f_\lambda(t/S),\qquad x=S v(\lambda).
\]
Then $F(x)=x$, the map $F$ is monotone, and
\begin{equation}
 F(t+c\one)=F(t)+\gamma c\one\qquad(t\in\R^n,\ c\in\R).
\label{eq:application-shift}
\end{equation}
Thus $F$ has the properties used in Section~\ref{sec:labeling}.
Define the integer vectors
\begin{equation}
 Y=\lfloor x\rfloor-N\one,
 \qquad Z=\lceil x\rceil+N\one.
\label{eq:comparison-vectors}
\end{equation}
Floors and ceilings are coordinatewise; in particular, a floor of a
negative number is taken toward minus infinity. These vectors are used
only in the proof. The algorithm will not compute them.

\begin{lemma}
\label{lem:comparison-vectors}
The vectors in \eqref{eq:comparison-vectors} satisfy
\[
 F(Y)\ge Y,\qquad F(Z)\le Z,\qquad Y\le x\le Z,
 \qquad \|Y\|_\infty,\|Z\|_\infty\le J.
\]
For every coordinate $i$,
\begin{equation}
 Z_i-Y_i\le 2(N+1)=K/128.
\label{eq:comparison-gap}
\end{equation}
\end{lemma}
\begin{proof}
The inequalities
$x-(N+1)\one\le Y\le x-N\one$ and
$x+N\one\le Z\le x+(N+1)\one$, together with
$\gamma(N+1)=N-1/N\le N$, give
\[
 F(Y)\ge x-\gamma(N+1)\one\ge x-N\one\ge Y
\]
and
\[
 F(Z)\le x+\gamma(N+1)\one\le x+N\one\le Z.
\]
Here we used monotonicity and \eqref{eq:application-shift}.
The remaining conclusions follow from these defining bounds and
$\|x\|_\infty\le SW$.
\end{proof}

\subsection{Shrinking the box}

Recall that the permitted widths are $D=h2^d$ with
$128\le h\le255$ and integers $d\ge0$. Start with
\begin{equation}
 A=-J\one,\qquad D=D_0=128\cdot2^{d_0},
 \label{eq:initial-box}
\end{equation}
where $d_0$ is the least nonnegative integer for which $D_0\ge2J$.
The box $[A,A+D\one]$ contains $Y$ and $Z$.

While $D>K$, call $\cL(A,D,a,b)$ with
$a_i=b_i=1/n$ for all $i$. Put $u=A+D\one$ and $P=D/32$.
Within the current box, $Y$ is a subsolution and $Z$ is a
supersolution, and their support masses are at most one. Consequently,
the labeling guarantee has the following contrapositive form:
\begin{equation}
 \begin{aligned}
 \text{label }+\text{ at }i&\quad\Longrightarrow\quad Z_i>A_i+P,\\
 \text{label }-\text{ at }i&\quad\Longrightarrow\quad Y_i<u_i-P.
 \end{aligned}
 \label{eq:label-exclusions}
\end{equation}
For example, $Z_i\le A_i+P$ would be a minus claim and would require
a minus label.

Let $D'$ be the next smaller permitted width. More explicitly,
\begin{equation}
 D'=\begin{cases}
 (h-1)2^d,&h>128,\\
 255\cdot2^{d-1},&h=128.
 \end{cases}
 \label{eq:next-width}
\end{equation}
The second case has $d\ge1$, because the loop condition implies
$D>K>255$. Set $\delta=D-D'$. In the first case
$\delta/D=1/h\le1/128$; in the second case this ratio is $1/256$.
Thus
\begin{equation}
 0<\delta\le D/128,
 \qquad Z_i-Y_i\le K/128<D/128.
 \label{eq:shrink-slack}
\end{equation}
Define
\[
 A'_i=\begin{cases}
 A_i+\delta,&\text{if the label is }+,\\
 A_i,&\text{if the label is }-.
 \end{cases}
\]
Replace $(A,D)$ by $(A',D')$.

\begin{lemma}
\label{lem:box-invariant}
Every box in this iteration contains $Y$ and $Z$. The iteration makes
at most $128(d_0+1)$ calls to $\cL$ and ends with $D\le K$.
\end{lemma}
\begin{proof}
The initial containment was established above. At a plus coordinate,
\eqref{eq:label-exclusions} and \eqref{eq:shrink-slack} give
\[
 Y_i>Z_i-D/128>A_i+D/32-D/128
       =A_i+3D/128\ge A_i+\delta.
\]
The lower endpoint can therefore increase by $\delta$; the upper
endpoint remains unchanged. At a minus coordinate the same inequalities
give
\[
 Z_i<Y_i+D/128<u_i-D/32+D/128
       =u_i-3D/128\le u_i-\delta.
\]
The upper endpoint can decrease by $\delta$; the lower endpoint remains
unchanged. This proves the induction step for both comparison vectors.

There are only $128$ permitted mantissas at any exponent, and the
exponent never increases. The strictly decreasing sequence of widths
therefore has at most $128(d_0+1)$ terms before reaching the stopping
range, which contains all permitted widths at exponent zero.
\end{proof}

At termination, $Y\le x\le Z$ and the invariant imply
$A_i\le x_i\le A_i+D$. Hence
\begin{equation}
 \left|\frac{A_i}{S}-v_i(\lambda)\right|
 \le\frac KS=\frac1{16H}.
 \label{eq:fixedpoint-error}
\end{equation}
Theorem~\ref{thm:discount} now gives
\begin{equation}
 \left|\frac{A_i}{S}-\val(i)\right|\le\frac1{8H}.
 \label{eq:final-error}
\end{equation}
Output exactly the indices satisfying
\begin{equation}
 \frac{A_i}{S}\ge-\frac1{2H},
 \quad\text{equivalently}\quad A_i\ge-8K.
 \label{eq:output-test}
\end{equation}
If $\val(i)=0$, \eqref{eq:final-error} puts $A_i/S$ at least
$-1/(8H)$, so the index is included. If $\val(i)>0$, the same conclusion
holds. If $\val(i)<0$, the value gap in Theorem~\ref{thm:discount}
gives $A_i/S\le-1/H+1/(8H)=-7/(8H)$, so the index is excluded.
This proves exact correctness, including equality at zero.

\subsection{Bit complexity and uniformity}
\label{subsec:bit-complexity}

It remains to show that the potentially large integer widths do not
produce a large recursion tree, and that every local computation has
polynomial bit cost. The first point uses the mass-product rank of
Section~\ref{sec:labeling}; the second uses the explicit binary encoding.

Since every vertex has an outgoing edge and the input is explicit,
$n,|E|=O(L+1)$. The binary lengths of the probabilities and rewards give
\[
 \log_2 Q=O(L+1),\qquad \log_2 W=O(L+1),\qquad
 \log_2 A_0=O((L+2)^2).
\]
In particular, all integers in \eqref{eq:application-parameters} and
\eqref{eq:initial-box} have $O((L+2)^2)$ bits, and
$d_0=O((L+2)^2)$. Reading the input and forming these integers by ordinary
exact arithmetic take polynomial time. Every recursive box lies in its
parent, and every outer box lies in the initial box. Thus all box
endpoints, pivot coordinates, and integer arguments to $F$ have
polynomial binary length.

For an integer vector $t$, the quantities used to evaluate $F_i(t)$ are
\begin{equation}
 \frac{S w(e)+(N-1)t_j}{N},\qquad e:i\longrightarrow j.
 \label{eq:exact-evaluation}
\end{equation}
Max and Min take exact extrema. At a chance vertex write the encoded
probability as $p(e)=c_e/d_e$ with $d_e>0$. Its weighted sum can be
formed with the common denominator $QN$ and integer numerator
\[
 \sum_{e:i\to j}c_e\frac{Q}{d_e}
           \bigl(Sw(e)+(N-1)t_j\bigr).
\]
Every $d_e$ divides $Q$, even when input fractions are not reduced.
This formula gives polynomial bit lengths and polynomial bit cost for
an evaluation and its coordinatewise floors. Zero probabilities and
parallel edges require no special handling.

Consider one outer call, whose initial masses are $a_i=b_i=1/n$.
For a recursive call put
\[
 \mu=\min_i a_i b_i,\qquad \rho=\frac{32}{25},\qquad
 B(\mu)=\min\{j\in\Z_{\ge0}:\rho^j\mu>1\}.
\]
The initial value $B_0=B(1/n^2)$ is $O(\log(n+1))$; this includes
$n=1$, for which $B_0=1$. Every recursive edge decreases $d+B$.
Only the provisional-label child can preserve $B$, and it decreases
$d$ by one. Each other child decreases $B$ by at least one and does
not increase $d$. Thus the depth is at most
\[
 \Lambda=d_0+B_0,
\]
and any root-to-node path contains at most $B_0$ edges decreasing $B$.
At each node there is at most one $B$-preserving child and at most
\[
 s_*=2(1024n+1)+1=2048n+3
\]
other children: the two pivot passes account for the first term and
the reweighted call for the final one.

A path of length $\ell$ with $j$ edges decreasing $B$ is specified by
their $j$ positions and at most $s_*$ choices at each such position;
the other child, if present, is unique. The total number of nodes is
therefore at most
\begin{align}
 \sum_{\ell=0}^{\Lambda}\sum_{j=0}^{\min(\ell,B_0)}
       \binom{\ell}{j}s_*^j
 &\le (\Lambda+1)(B_0+1)(1+\Lambda s_*)^{B_0} \\
 &=2^{O((\log_2(L+2))^2)}.
 \label{eq:tree-bound}
\end{align}
The distinction between these two kinds of children is essential:
the depth may be polynomial in $L$, but branching choices occur only
$O(\log(n+1))$ times on any path.

All masses can be stored as unreduced positive fractions. An edge that
changes masses multiplies each by one of a fixed set of rational
constants, and so increases its numerator and denominator lengths by
at most a constant. The other edge leaves masses unchanged. Thus the
mass lengths are polynomial (indeed $O(\log(n+1))$), and every comparison,
update, and support-independent base test has polynomial bit cost.
The algorithm does not compute $B$; it is used only for this analysis.

Outside its children, a nonbase call performs $O(n+1)$ pivot updates
and associated vector operations, all of polynomial bit cost. In a
width base case, $d<12$ gives $D\le255\cdot2^{11}$, an absolute
constant, so the clipped-floor iteration uses at most $1+nD$
polynomial-cost evaluations. The mass base case is also polynomial.
Consequently \eqref{eq:tree-bound} bounds the bit cost of one outer
call up to a polynomial factor. Lemma~\ref{lem:box-invariant} bounds
the number of outer calls by a polynomial as well.

Finally, all loops, tie rules, integer operations, and recursive calls
are specified by finite algorithms independent of the input. A
depth-first implementation stores a stack of polynomial depth whose
individual records have polynomial length. Scanning such a stack or
copying a record on a deterministic Turing machine costs at most a
polynomial factor per local operation. The final integer tests
\eqref{eq:output-test} and the output take polynomial time.
Absorbing these factors into \eqref{eq:tree-bound} proves the bound
$2^{C(\log_2(L+2))^2}$ for one absolute constant $C$. This completes
the proof of Theorem~\ref{thm:main}.
